\section*{Chapter \ref{ch:s1:calendar-spreads-and-curve-trades} --- Calendar Spreads and Curve Trades}

\begin{solution}{exo:s1:calendar-spreads-and-curve-trades:1}
$\ln 0.5/\ln 0.949 = 13.2$ trading days (13.3 from the unrounded coefficient).
\end{solution}

\begin{solution}{exo:s1:calendar-spreads-and-curve-trades:2}
$-0.01 \times 252/21 = -0.12$, a contango of 12\% a year. A carry of $-575.5\%$ is a log spread of $-5.755 \times 21/252 = -0.48$, so the near
price was $e^{-0.48} = 0.62$ times the far one.
\end{solution}

\begin{solution}{exo:s1:calendar-spreads-and-curve-trades:3}
On that day contract~1 is a different contract (yesterday's contract~2): the change in ``contract 1 minus contract 2'' mixes the old and new pairs, and
the jump between them is not a P\&L anyone earned. The spread must be followed by contract identity.
\end{solution}

\begin{solution}{exo:s1:calendar-spreads-and-curve-trades:4}
About $0.30 \times 10\% - 0.05 \times 9.9\% = 2.5\%$ a year, most of the gross return.
\end{solution}

\begin{solution}{exo:s1:calendar-spreads-and-curve-trades:5}
In March contango widened far beyond anything in the rule's twenty-day window: the rule was long the spread (betting on narrowing) while it kept widening
as storage filled. In April it re-entered with smaller positions, sized on the much higher volatility, and part of the spread's move reversed.
\end{solution}

\begin{solution}{exo:s1:calendar-spreads-and-curve-trades:6}
It was chosen among four thresholds after seeing the one storage crisis that decides the result; any threshold between normal contango and the 2020
levels would look good in hindsight. Evidence would need a threshold fixed in advance from the mechanism (storage costs, inventories) and tested on
crises not used to choose it.
\end{solution}

\begin{solution}{exo:s1:calendar-spreads-and-curve-trades:7}
Daily, the z-score's small day-to-day changes move the position every day, and at 4 basis points per leg per unit traded the turnover costs more than the
slow signal earns ($-0.29$ after costs against 1.00 before). Monthly, the position changes twelve times a year; the signal, whose information decays
over a year, loses little by waiting, and 0.76 of 1.05 survives.
\end{solution}

\begin{solution}{exo:s1:calendar-spreads-and-curve-trades:8}
The legs hedge the level, not the slope. When storage runs out the slope has no bound, and at ten times leverage a move like March 2020 (58\% of capital
at a 10\% volatility target) would be a multiple of the capital. Twenty years without a storage crisis is not evidence of none.
\end{solution}

\begin{solution}{pb:s1:calendar-spreads-and-curve-trades:1}
\begin{enumerate}
\item Positions in several deliveries neutral to the level, trading the shape; the return of a spread toward its normal level.
\item $d(\ln F(T_a) - \ln F(T_b)) = (T_b - T_a)\,dc$, since the level and seasonal terms cancel.
\item Convenience yields fall nonlinearly as inventories rise; the basis and past returns reflect inventories and forecast premiums.
\item The cost of storage (full carry); it fails when storage is full.
\item By serial number: the pair clear of expiry is held and followed from one day to the next.
\item A daily autocorrelation of 0.949, a half-life of 13.3 days.
\item A 20-day z-score, entry beyond 1.5, exit inside 0.5 or on a sign change, sized to 10\% volatility, one cent a barrel per unit traded.
\item 0.30 and 0.05.
\item The May contract settled at $-\$37.63$, the first negative price in 37 years, amid oversupply, a demand collapse and storage at Cushing near
capacity.
\item The risk that contango widens without bound when storage runs out.
\item Long the spread from 10 January, carry falling to $-216\%$ by 31 March and $-575\%$ on 21 April; $-58.3\%$ in March, $+22.9\%$ in April.
\item Losses of 28.1\%, 17.8\% and 8.8\% for thresholds of $-80\%$, $-40\%$ and $-20\%$; the best was chosen after the fact.
\item \textbf{Named result.} From 1986 to 2019 the WTI spread mean-reversion rule had a Sharpe ratio of 0.30 before costs and 0.05 after; in 2020 it lost
58.3\% of capital in March, made back 22.9\% in April, and lost 42.7\% over the year at a 10\% volatility target.
\item Inventories and storage costs, fixed before the test.
\item A slow signal must be traded slowly: $-0.29$ daily, 0.76 monthly after costs.
\item The basis did not forecast spot VIX but did forecast VIX futures returns; the hedged contango trade was profitable.
\item A volatility spike that turns contango into backwardation.
\item The roll-period calendar spread.
\item Chapter~20 traded the level of carry across markets; a calendar spread trades its change within one market.
\item It hedges the curve's level, not its slope, and the slope is unbounded when storage runs out.
\end{enumerate}
\end{solution}

\begin{solution}{iq:s1:calendar-spreads-and-curve-trades:1}
Storage costs and interest (the upper bound on contango), the convenience yield of holding inventory (which rises as inventories fall, giving
backwardation), seasonality of supply and demand, and risk premiums paid by hedgers.
\end{solution}

\begin{solution}{iq:s1:calendar-spreads-and-curve-trades:2}
One series per contract (by delivery month), with its expiry and first notice dates, plus a roll calendar; generic series are derived views, never the
stored truth, and every P\&L is computed contract by contract.
\end{solution}

\begin{solution}{iq:s1:calendar-spreads-and-curve-trades:3}
Check inventories and the cost of storage at the delivery point, the positions of \omterm{def:s1:index-rebalancing:fund}{index funds} and ETFs near their roll, and the exchange's notices; cut
the position if the contango is approaching or exceeding full carry, rather than adding to it on a mean-reversion signal.
\end{solution}

\begin{solution}{iq:s1:calendar-spreads-and-curve-trades:4}
Exchanges usually margin spreads less than outrights because the legs offset; the book should be stressed on the slope itself, with historical storage
crises and hypothetical contango at and beyond full carry, and with the legs' liquidity near expiry.
\end{solution}

\begin{solution}{iq:s1:calendar-spreads-and-curve-trades:5}
Buyers of protection pay a premium, so futures sit above the expected VIX and roll down toward spot; the risk is a spike, when spot VIX jumps above the
futures and the short loses many months of premium in days.
\end{solution}

\begin{solution}{iq:s1:calendar-spreads-and-curve-trades:6}
With $s_{t+1} = \phi s_t + \varepsilon_{t+1}$, the P\&L is $-s_t(s_{t+1} - s_t) = (1 - \phi)s_t^2 - s_t\varepsilon_{t+1}$, with mean
$(1 - \phi)\sigma^2/(1 - \phi^2) = \sigma^2/(1 + \phi)$. Its variance is dominated by $E[s_t^2]\sigma^2 = \sigma^4/(1 - \phi^2)$ for $\phi$ near one, so
the daily Sharpe ratio is about $\sqrt{(1 - \phi)/(1 + \phi)}$. The position changes by $s_{t+1} - s_t$, about $\varepsilon$, so costs take about $k\,
\sigma\sqrt{2/\pi}$ a day: the trade needs $\sigma/(1 + \phi) > k\sqrt{2/\pi}$, a spread noisy enough to pay for its own trading.
\end{solution}
